\begin{result}{E.2.12}{Finite equations give a comparison in an admissibly generated integral category}{Proved with explicit lattice hypotheses}
Let $R$ be noetherian, let $\mathcal C$ be a small rigid symmetric
$R$-linear exact tensor category, and let
\[
 \omega_1,\omega_2:\mathcal C\longrightarrow\operatorname{Proj}_f(R)
\]
be faithful exact tensor functors. Assume every object is obtained,
up to isomorphism, from one object $M$ by finite sums, tensor products,
duals, and admissible subobjects and quotients, and that
$\omega_1(M),\omega_2(M)$ are finite free of the same rank.
Then $\Isom^\otimes(\omega_1,\omega_2)$ is represented by a closed
subscheme of $\Isom_R(\omega_1(M),\omega_2(M))$. A finite diagram of
original objects, maps and admissible sequences determines each
comparison uniquely.

For $R=\Z$, consider such a category of abelian motives over the
quartic field of E.2.6 with its two integral Betti realizations.
Require that the generator's lattices are given by a specified
admissible expression from the actual Betti lattices of finitely
many abelian and Artin generators. Then these two realizations
are integrally tensor isomorphic on $\mathcal C$. In this
specified integral model, a single admissible tensor generator
does not distinguish those quartic places.
\end{result}
\begin{proof}
After bases, the ambient isomorphism scheme has coordinate ring
$B=R[u_{ij},\det(u)^{-1}]$, with universal matrix $u$ and inverse
$v$. Finite sums, tensor products and duals give the corresponding
maps $u_P,v_P$. For later objects whose realizations are projective
but not free, realize each as an idempotent summand of a finite
free module and write the equations there. Equivalently, evaluate
each module-valued equation against the dual of its finite
projective Hom module to obtain its coefficient ideal. The
equations are therefore global over $R$, without bases on every
projective module.

For an admissible monomorphism $i:N\to P$, with $P$ already
constructed, put $i_a=\omega_a(i)$ and choose retractions $r_a$
with $r_ai_a=1$. They exist because the cokernel is projective.
Impose
\[
 (1-i_2r_2)u_Pi_1=0,\qquad
 (1-i_1r_1)v_Pi_2=0,
 \tag{E.2.12a}
\]
and define
\[
 u_N=r_2u_Pi_1,\qquad v_N=r_1v_Pi_2.
 \tag{E.2.12b}
\]
The equations preserve the first image under $u_P$ and the
second under $v_P$. Split identities persist after extension to
every $R$-algebra $S$. Substitution therefore gives
$u_Nv_N=1=v_Nu_N$ over every $S$, including nonflat base changes.
For an admissible quotient $q:P\to N$, preserve its kernel in
both directions and choose sections $s_a$ of $q_a=\omega_a(q)$.
The induced inverse quotient maps are
\[
 u_N=q_2u_Ps_1,\qquad v_N=q_1v_Ps_2.
 \tag{E.2.12c}
\]
They are independent of the sections because the differences
lie in the preserved kernels; their composites are identities
because $s_aq_a-1$ has image in those kernels.

Choose one finite admissible expression for every object, apply
these formulas, and impose every original naturality equation
\[
 u_Y\omega_1(f)=\omega_2(f)u_X\qquad(f:X\to Y),
 \tag{E.2.12d}
\]
together with the tensor, unit, symmetry and duality equations.
The original isomorphisms between presentations are included, so
the resulting maps are independent of the chosen expressions.
Every solution is a tensor natural isomorphism, and conversely
every such isomorphism satisfies the equations. Its component on
$M$ determines all the other components. This proves the stated
closed-subscheme representation.

Let $I\subset B$ be the ideal of all coefficient equations.
Noetherianness makes $I$ generated by a finite subset of the
\emph{original} equations: choose finitely many ideal generators,
express each using finitely many original coefficients, and take
their union. Include the finite admissible presentations of the
involved objects. This gives the finite original diagram.
Prescribed components on a fixed Artin or CM object add ordinary
matrix equations and are covered by the same argument.

For the geometric application, use the element
$\sigma\in\Gal(\Qb/F_{\rm CM})$ of E.2.6 exchanging the specified
real roots. The finite-diagram construction (E.2.11c), before its
real Hodge correction, gives an integral map $T_0$ on the finite
defining diagram. It preserves all the requested rational absolute
Hodge tensors and has the $\sigma$-selected Artin permutation.
The morphisms have integral Betti matrices because they are
morphisms of $\mathcal C$. The admissible expressions recover
saturated image lattices and free quotient lattices: an integral
ambient isomorphism preserving a rational image subspace preserves
its intersection with the ambient lattice, and then induces the
isomorphism on the free quotient. Thus $T_0$ satisfies the finite
original equations. Its component on $M$ satisfies all of $I$,
and (E.2.12a)--(E.2.12d) extend it uniquely to the asserted
global integral comparison.

The lattice hypothesis is part of the assertion. The injection
on $H_1$ induced by a nontrivial isogeny $q:E\to E/H$ has finite
torsion cokernel, so that injection is not an admissible subobject
construction in the exact category of finite free $\Z$-modules.
This does not exclude $q$ as an ordinary morphism when both
objects have independently supplied admissible expressions:
then (E.2.12d) includes its naturality. Nor does it assert that
every isogenous target lacks such an expression. It states the
additional presentation required before applying the result to
that lattice. The entire geometric diagram in E.2.9 has not been
identified with this admissibly generated category, and the full
pure integral assertion (E.2.3d) is not concluded.
\end{proof}
